%% kapa: exam=1 round=24 year=2013 period=1 form=A subject=civil_law subject_ko=민법 seq=39 number=39
%% source: Q-Net 감정평가사 기출문제 (cst00309) artlSeq=5073802 file=extracted/1008/5073802/2013년 제24회 감정평가사 1차시험 기출문제/1교시(A형).hwp
%% conversion: kordoc 4.24.51 -> tools/gen1.py (auto); review=proofread-sonnet; flags=none
\begin{question}{39}
저당권의 효력에 관한 설명으로 옳지 않은 것은? (다툼이 있으면 판례에 의함)
\begin{choices}
  \item 토지저당권의 효력은 토지상에 타인이 무단 경작하여 성숙한 농작물에는 미치나, 적법한 토지임차인이 경작한 농작물에는 미치지 않는다.
  \item 건물에 저당권이 설정된 경우 그 건물에 연이어 설치되어 그 일부에 불과한 부속건물에도 저당권의 효력이 미치므로 등기부에 등재되지 않은 그 부속건물을 경매목적에 포함시킨 매각허가 결정은 적법하다.
  \item 저당권의 효력은 저당부동산에 대한 압류가 있은 후에 저당권설정자가 그 부동산으로부터 수취한 과실 또는 수취할 수 있는 과실에 미친다.
  \item 구분건물의 전유부분에 대한 저당권의 효력은 특별한 사정이 없으면 대지사용권에도 미친다.
  \item 저당권자가 물상대위권을 행사하기 위한 요건인 ‘지급 또는 인도 전 압류’는 반드시 저당권자가 할 필요는 없다.
\end{choices}
\end{question}
\answer{1}
